% Rupture du courant. Résolution analytique
Une bobine d'inductance $L$ et de résistance $r$, placée en série avec un
conducteur ohmique de résistance $R^{\prime}$ est soumise à une tension qui
passe d'une valeur $E=\qty{6.0}{\V}$ à une valeur nulle. On note $R$ la
résistance totale du circuit.

\begin{enumerate}
    \item Faire un schéma du circuit.
    \item En appliquant la loi d'additivité des tensions, établir une
          relation~(1) entre $u_{R}$, tension aux bornes du conducteur ohmique
          et $u_{L}$ tension aux bornes de la bobine.
    \item En exprimant $u_{R^{\prime}}$ et $u_{L}$ en fonction de $i$, établir
          l'équation différentielle à laquelle obéit $i$.
    \item Une solution de cette équation différentielle est de la forme~:
          $i=c+a\times\mathrm{e}^{bt}$.
          \begin{enumerate}
              \item Déterminer les constantes $a$, $b$ et $c$.
              \item En déduire l'expression de l'intensité $i$ en fonction de la
                    date $t$.
          \end{enumerate}
\end{enumerate}


